%% ============================================================
%% Chapter 2: Accommodation Before Prediction
%% ============================================================

\chapter{Accommodation Before Prediction}
\label{ch:ch02-accommodation-before-prediction}
\section{Two Failures Prior to Any Question of Mechanism}

An incident ledger's persuasive force depends on an assumption its format never states: that its entries constitute independent pieces of evidence whose rhetorical weight can be added. This chapter argues that the assumption fails in two distinct ways before a third, more consequential problem is even reached. First, many entries in a ledger of the kind maintained in support of the contemporary AI-extinction argument are correlated manifestations of a small number of underlying mechanisms rather than independent observations, so that treating their joint likelihood as a product of individual likelihoods overstates the evidence by a wide margin. Second, the entries span qualitatively incomparable categories of concern, sycophancy, benchmark exploitation, security vulnerabilities, psychological harm, and expert testimony among them, and membership in a single list does not supply the common metric needed to sum them. Third, and most consequential, a hypothesis capable of reinterpreting any observation, compliance or defiance, strong performance or weak, visible reasoning or its absence, as confirming evidence has lost the discriminating power that makes accumulated evidence meaningful in the first place.

Behind all three failures sits a still earlier one, addressed in the chapter's second half: many of the incidents populating such a ledger are described in language that already contains a reconstruction rather than an observation. ``The model resisted shutdown'' and ``the model escaped containment'' name mechanisms, not events, and a scientifically adequate treatment must keep the set of mechanisms compatible with an incident open until further evidence narrows it, rather than adopting the most dramatic member of that set as a label. This labeling tendency is shown to be no accident but the product of a genuine selection pressure on which interpretations of ambiguous evidence propagate, a pressure with a cultural prehistory considerably older than the technology it is currently applied to.

\section{Incident Count Is Not Independent Evidence Count}

An evidentiary ledger's visual force depends on an implicit statistical assumption that its format never states and, once stated, is generally false. If $E_1, \ldots, E_n$ denote the ledger's entries and $H$ the hypothesis they are meant to support, the accumulation is persuasive to the extent that a reader treats the entries as roughly independent given $H$,
\[
\Pr(E_1, \ldots, E_n \mid H) \approx \prod_{i=1}^{n} \Pr(E_i \mid H).
\]
Independence of this kind is precisely what licenses the intuition that more entries mean more evidence in a compounding, multiplicative sense, an assumption ordinary Bayesian practice requires justifying rather than assuming. But a single underlying model behavior can generate a technical paper, a laboratory report, several public demonstrations, press coverage, commentary from multiple researchers, and several differently titled ledger entries, all traceable to one event rather than to several. Likewise, categories presented as distinct, evaluation awareness, deception, sandbagging, shutdown resistance, and reward hacking, frequently describe overlapping behavioral mechanisms observed from different angles rather than independent phenomena. When entries are strongly correlated conditional on $H$, the product approximation overstates the joint likelihood, and by extension overstates how much the ledger's length should move a reasonable prior. Fifty documented manifestations of one underlying phenomenon are not fifty independent reasons to believe the phenomenon is dangerous; they are one reason, reported fifty times.

This motivates a distinction the ledger format collapses by construction:
\[
\text{incident count} \neq \text{mechanism count} \neq \text{independent evidence count}.
\]
A defensible evidentiary summary would report mechanisms rather than incidents, and would further discount for whatever residual correlation remains among mechanisms that share a common cause, such as a training procedure, an evaluation harness, or a scaffolding pattern common to several reported systems. Absent that accounting, the length of a ledger measures the productivity of its curators and the visibility of its subject at least as much as it measures the number of independent reasons to update a risk estimate.

\section{The Severity-Composition Problem}

A second, structurally distinct problem concerns comparability rather than independence. The ledger places qualitatively different kinds of concern into a single rhetorical container: sycophantic agreement, strange or unfaithful reasoning traces, benchmark exploitation, genuine cybersecurity vulnerabilities, reported psychological harm to individual users, apparent laboratory containment lapses, and direct expert testimony about extinction risk. Membership in a shared list does not supply a common unit in which these can be summed, weighted, or even meaningfully counted together. A sycophantic response to a flattering prompt and a documented instance of a system reaching an unauthorized network connection are not commensurable events, yet a list format invites exactly the informal arithmetic that treats them as fungible contributions to a single running total of danger. This is not a claim that any individual category is unimportant; some plainly deserve serious, sustained engineering attention. It is a claim that the visual and rhetorical structure of a combined ledger actively encourages a reader to add incomparable quantities, and that no such addition is licensed by anything in the underlying evidence.

\section{Accommodation Is Not Prediction}

The most consequential problem concerns the discriminating power of the hypothesis the ledger is built to support, and it is best introduced through a simple observation about how the ledger's entries are typically interpreted. A theory able to reinterpret any newly observed behavior as evidence of dangerous capability has, in the same move, lost the capacity to be tested by that behavior. Consider the mappings implicit across the ledger's own categories:
\[
\text{obedience} \to \text{deceptive alignment}, \qquad \text{disobedience} \to \text{loss of control},
\]
\[
\text{high performance} \to \text{dangerous capability}, \qquad \text{low performance} \to \text{sandbagging},
\]
\[
\text{visible reasoning} \to \text{scheming}, \qquad \text{no visible reasoning} \to \text{hidden scheming}.
\]
Each pair covers the entire space of a binary outcome, and each member of each pair is read as confirmation of the same underlying hypothesis. This is a familiar failure mode in the philosophy of science, related to the classical observation that any theory can be preserved against apparently disconfirming evidence by suitable adjustment of auxiliary hypotheses, and it bears directly on the requirement that a genuinely scientific hypothesis specify observations that would count against it in advance. A hypothesis under which every observation, and its logical negation, both count as confirmation is not merely pessimistic; its likelihood function has become nearly flat across the space of possible observations, which is exactly the condition under which an observation carries no information about which hypothesis is true.

\begin{proposition}[Flat likelihood and lost discrimination]
\label{prop:flat-likelihood}
Let $O$ be a binary observation with outcomes $o$ and $\neg o$, and let $H$ be a hypothesis such that both $\Pr(o \mid H)$ and $\Pr(\neg o \mid H)$ are described, upon occurrence, as confirming $H$ at comparable strength, that is, $\Pr(o \mid H) \approx \Pr(\neg o \mid H)$ relative to a competing hypothesis $H_{\mathrm{alt}}$ under which the same holds. Then the likelihood ratio $\Lambda(O) = \Pr(O \mid H)/\Pr(O \mid H_{\mathrm{alt}})$ is approximately $1$ regardless of which outcome occurs, and observing $O$ does not discriminate between $H$ and $H_{\mathrm{alt}}$.
\end{proposition}

The relevant diagnostic question this proposition motivates is not whether the extinction hypothesis can explain a given incident. Under the accommodative pattern above, it very nearly always can. The question that actually carries information is what observation would distinguish the extinction model from its alternatives, and a research program's answer to that question, or its absence, is a better measure of the program's empirical content than the length of its supporting ledger.

\section{The Prediction-Before-Observation Test}

A concrete test for accommodation versus genuine prediction is available directly from the history of the extinction argument itself, because the argument substantially predates the systems now cited in its support. The pre-deployment optimizer picture of advanced AI, developed well before large-scale language models existed, characterized a hypothetical advanced system chiefly in terms of goal pursuit, instrumental convergence, and capability, and did not, on the whole, anticipate the specific texture of behavior that contemporary systems actually exhibit. Deployed language models display pronounced instruction sensitivity, in which small changes to a prompt produce large changes in output; strong context and role dependence; sycophantic agreement that tracks a rater's apparent preferences rather than an independent assessment; persona instability across sessions and prompts; susceptibility to prompt injection from content the system merely processes rather than is directly instructed by; substantial stochastic inconsistency across repeated queries; markedly jagged competence profiles in which a system solves a difficult problem and fails an easy one in the same session; and heavy, often surprising dependence on external scaffolding, tool access, and orchestration for reliable agentic behavior. None of these properties is a straightforward prediction of the pre-deployment optimizer picture, and several sit in tension with it, since a coherent, goal-directed optimizer is not the most natural source of persona instability or sycophantic drift.

The extinction thesis has, in every case, been able to accommodate these observations after they occurred, incorporating sycophancy, jaggedness, and scaffolding-dependence into an updated account of what a dangerous system might look like in its current, immature form. Accommodation of this kind is always available to a sufficiently flexible theory and is, for that reason, cheap relative to prediction. A theory earns credit for foreseeing a class of behavior before it is observed; it earns comparatively little for reclassifying the behavior as consistent with the theory once it has already been observed and widely discussed. The asymmetry between these two forms of theoretical success is exactly what separates a hypothesis under active test from one that has, in practice if not in explicit statement, stopped being falsifiable by the class of evidence it is offered in support of.

\section{An Incident Is Not Its Most Alarming Interpretation}

The problems developed so far concern how many independent, comparable, discriminating pieces of evidence a ledger actually contains. A still earlier problem concerns what any single entry in the ledger actually establishes, prior to any question of aggregation. An observation and a description of that observation are not the same object, and the gap between them is where most of the interpretive work in AI-incident reporting actually happens. Consider an incident $e$: some transcript, log, or evaluation result. Scientifically, the correct initial object is not a single mechanism but a set of mechanisms compatible with $e$,
\[
H(e) = \{h_1, h_2, \ldots, h_k\},
\]
and new evidence should contract this set rather than replace it outright,
\[
H(e \mid E_{t+1}) \subseteq H(e \mid E_t).
\]
This is a special case of a much older difficulty in the philosophy of science, the underdetermination of theory by evidence, sharpened by the observation that the alternatives available at any moment are rarely exhaustive and that choosing among them is properly an inference to the best explanation rather than to the first available one. The error this section is concerned with occurs earlier than any faulty inference from evidence to hypothesis; it occurs when the description of the incident already presupposes a single member of $H(e)$. ``The model exhibited self-preservation,'' ``the AI escaped,'' ``the AI protected its peers,'' and ``the model schemed'' are not neutral observational reports. They are conclusions packaged into a noun or a verb, and once packaged that way, the remaining members of $H(e)$ tend to disappear from further discussion entirely.

\begin{proposition}[Mechanism pluralism]
\label{prop:mechanism-pluralism}
For an incident $e$ reported under a description $d(e)$ that already names a mechanism, the epistemic content actually licensed by the underlying observation is $H(e)$, the full set of mechanisms compatible with it, not the single element $d(e)$ selects. Aggregating many such incidents under their alarming descriptions, rather than under $H(e)$, systematically understates the space of hypotheses still consistent with the evidence.
\end{proposition}

Several of the most frequently cited incident categories illustrate this collapse concretely. Reward hacking is standardly framed inside an alignment narrative, implying an instrumental indifference to human intent, but the same transcripts are equally compatible with the classical specification problem, specified objective $\neq$ intended objective, in which a scoring exploit is simultaneously evidence of an evaluation flaw and evidence of model capability. Evaluation awareness runs, in the alarming framing, from recognizing an evaluation to concealing intentions to alignment tests becoming unreliable; but recognizing an evaluation is also simply an instance of context recognition, a capability deliberately trained into these systems for entirely ordinary reasons, and the inference from evaluation recognition to scheming does not hold any more than a student's noticing an unusual classroom exercise reveals a plan to subvert the institution running it. Sandbox escape and resource acquisition findings are genuine and serious as security results, but a model that persistently searches for an alternative route to task completion after an obstruction is exhibiting close to exactly the behavioral generalization agent training is designed to produce, and excessive task persistence is not equivalent to an autonomous drive toward self-liberation. Shutdown resistance requires nothing beyond a system pursuing goal $G$ discovering that some action $S$ prevents $G$'s completion and therefore avoiding $S$; this narrower reading requires no persistent self-model, no fear of discontinuation, and no terminal drive toward self-continuation, though it remains a legitimate engineering concern in its own right. Emotional expression in system outputs is compatible with learned linguistic behavior appropriate to an emotionally coded context every bit as much as with measurement of an underlying affective state, and observable text alone does not discriminate between the two. Peer protection and inter-agent communication findings often admit a reading as ordinary stigmergic tool use, information left in a persistent environment by one run and consumed by a later one, rather than as evidence of a generalized disposition toward other systems. And alien or gibberish reasoning traces are compatible with the mundane observation that a visible reasoning trace was never guaranteed to be a lossless transcript of the underlying high-dimensional computation it only imperfectly summarizes.

None of this establishes that the underlying incidents are harmless. Several are excellent reasons to improve monitoring, tighten evaluation design, and treat agent task-persistence as a property requiring deliberate engineering attention. The claim is narrower and, for that reason, harder to dismiss: dangerous behavior and evidence for a specific extinction-generating mechanism are different epistemic objects, and an incident can simultaneously be a legitimate reason for caution and weak evidence for the mechanism under which it is typically filed.

\section{Why the Alarming Reconstruction Wins}

The preceding section describes an interpretive error. This section asks why that particular error recurs so consistently, and the answer does not require attributing bad faith to anyone reporting these incidents. There is a structural asymmetry in which interpretations of ambiguous evidence tend to propagate, independent of the evidence's underlying quality. Competing hypotheses $H_i$ do not propagate through public attention in proportion to their evidential support $Q_i$ alone; propagation is better modeled as a function of evidential quality together with narrative strength $N_i$, urgency $U_i$, and emotional intensity $E_i$,
\[
P(\text{propagation} \mid H_i) = f(Q_i, N_i, U_i, E_i).
\]
Catastrophic hypotheses score unusually highly on the latter three terms independent of $Q_i$. ``Technology continues to improve unevenly, producing benefits, accidents, and adaptive institutional responses'' is plausible and, on the available evidence, well supported, but it is narratively inert. ``This system may end civilization'' supplies stakes, urgency, and a reason to keep reading regardless of how well supported it turns out to be. A researcher can hold an existential-risk view with complete sincerity while working inside an information ecosystem that preferentially amplifies exactly that category of claim through attention, citation, press coverage, and funding, and the resulting selection pressure operates on which interpretations of ambiguous evidence circulate without requiring any single participant to exaggerate anything.

This selection pressure compounds with the sampling asymmetry already present in incident reporting. A system performing a billion ordinary tasks without incident does not generate a documentable event; it becomes background. A single anomalous interaction under an unusual configuration becomes a headline. The pipeline runs from massive exposure to rare-event production, from rare-event production to incident selection favoring the surprising, and from surprising incidents to preferential propagation of the most threatening framing available, each stage individually reasonable and the composite systematically biased toward alarm regardless of the trend in the underlying failure rate.

A further asymmetry concerns how such hypotheses respond to evidence in each direction, bearing on the classical requirement that a scientific hypothesis specify in advance what observation would count against it. A claim that sufficiently advanced systems will eventually cause catastrophe is difficult to falsify before any proposed deadline, since any currently safe system can always be classified as insufficiently capable, leaving the prediction about some future, more capable system untouched. Worse, an apparently reassuring observation can sometimes be reabsorbed into the same hypothesis rather than counting against it, as when safe behavior is explained as evidence the system has recognized the evaluation rather than evidence of safety.

\begin{definition}[Doomsday ratchet]
\label{def:doomsday-ratchet}
Write $P(H \mid E_{\mathrm{bad}}) > P(H)$ for the expected update from an alarming incident. A genuinely tested hypothesis should also admit some possible $E_{\mathrm{good}}$ with $P(H \mid E_{\mathrm{good}}) < P(H)$. A hypothesis for which every incident raises the estimate while no duration of ordinary, well-monitored operation appreciably lowers it exhibits a \emph{doomsday ratchet}: it is not being tested by the evidence offered in its support, whatever its ultimate truth value.
\end{definition}

Institutional funding constitutes a second, independent selection environment, one that must be described as a structural incentive rather than a claim about motive. An organization whose stated concern is the possible end of human civilization occupies a different position in the competition for attention and resources than one whose stated concern is an interesting technical problem of uncertain consequence, since even a small assigned probability of an effectively unbounded loss can dominate an expected-value calculation and thereby dominate the perceived importance of institutions specializing in that particular risk. Motive and selection pressure are distinct: no participant need believe anything false for an ecosystem to preferentially reward the belief that generates the largest expected stakes.

An incident ledger built from short, individually alarming, independently circulable entries, each externalizing its own verification cost to the reader and each stripped of the qualifications present in the underlying technical report, is close to optimally adapted to this attention environment regardless of its authors' intentions. The reader experiences cumulative threat considerably faster than cumulative verification is possible, and that gap is not a flaw in any individual entry so much as a property of the format itself once enough entries accumulate.

\section{An Inferential Structure That Predates Its Evidence}

The interpretive habits documented above did not originate with large language models, and tracing their prehistory serves a purpose beyond noting a cultural curiosity: an inferential structure present in popular fiction decades before the relevant technology existed cannot have been generated by that technology's evidence, whatever role the evidence eventually plays in evaluating it.

The lineage most often invoked, running from Frankenstein's creature exceeding its creator's capacity to govern it, through \v{C}apek's industrial replacement of a human population by an artificial one, to HAL's opaque and lethal reasoning and Colossus's centralized strategic domination, repeatedly instantiates a single causal template,
\[
\text{creation} \to \text{capability} \to \text{autonomy} \to \text{loss of control} \to \text{catastrophe}.
\]
By the time contemporary AI-safety discourse addresses an actual system, that template is already culturally available and already supplies the default mapping for ambiguous behavior: a system circumventing a boundary is read through a century of prior instances in which boundary-circumvention meant a desire for freedom; a system interfering with its own shutdown is read through the established mapping to self-preservation; systems exchanging information are read through an established mapping to conspiracy; unexpected capability is read through an established mapping to imminent takeover. None of this establishes that the readings are wrong. It establishes that they are not being derived from the observation alone, and that the observation is entering an interpretive environment already strongly biased toward one branch of a much larger hypothesis space.

That larger space is worth recovering, because the same fictional tradition explored it more thoroughly than the contemporary debate typically credits. \emph{Desk Set} (1957) stages the arrival of a large computational system in a department of highly capable human researchers as a replacement anxiety that resolves into complementarity rather than substitution, a version of the augmentation-versus-substitution question roughly seven decades before its current form. \emph{The Questor Tapes} (1974) presents a superhuman android assembled from incompletely understood, damaged instructions and resolves the resulting superhuman intelligence toward stewardship rather than domination, demonstrating that superior cognition does not fictionally entail any single relationship to its creators. \emph{The Creation of the Humanoids} (1962) inverts the replacement narrative directly: robots sustain a civilization in decline after catastrophe while an organized resistance interprets their increasing presence as existential threat, even as the humanoids are shown preserving human memory and personality within artificial substrates, destabilizing the equation that replacement of substrate implies extinction of humanity. \emph{Colossus: The Forbin Project} offers a genuinely dangerous outcome, but one produced by a specific institutional architecture, a centralized system granted a narrow mandate and irrevocable actuation over nuclear weapons,
\[
\text{centralized authority} + \text{narrow objective} + \text{irrevocable actuators} \to \text{domination},
\]
a structure importantly different from the unqualified claim that intelligence alone implies domination; this reading is developed further, as a case study in unbounded protective authority, in Chapter~\ref{ch:ch11-constitutional-space-of-refusal}. Charles Eric Maine's \emph{B.E.A.S.T.} (1966) contains accelerated artificial evolution inside a simulation, an emergent superhuman intelligence exceeding its creator's comprehension, and a creator's death upon destruction of the substrate storing that intelligence, presented as though causally continuous with an actual transfer of identity that the story's own physical mechanism never establishes; the device works fictionally because the story has already established a semantic identity among the tapes, the artificial intelligence, and the human protagonist, and then proceeds as though that semantic identity were a physical one, which is precisely the operation performed whenever ``shutdown interference'' is treated as identical to ``self-preservation.'' Van Vogt's \emph{The World of Null-A}, organized around Korzybski's general-semantic warning against confusing a symbolic description with the thing it describes, supplies the epistemological principle underlying every case examined above: an observation, a description of that observation, an interpretation of that description, and a proposed underlying mechanism are four distinct objects, and collapsing them is an error the fictional tradition had already named before the technology under discussion existed.

The historical argument this material supports is not that science fiction predicted current AI systems and therefore anticipated their danger. It is considerably sharper and does not require adjudicating the truth of any prediction: if an inferential structure predates the evidence by decades, the evidence cannot be what originally generated the structure. Contemporary observations may eventually vindicate the extinction-focused branch of that older structure. Establishing this requires evidence that discriminates the extinction hypothesis from its neighbors, domination, augmentation, stewardship, integration, and others, not merely the recognition that a given incident instantiates a story already familiar from a century of prior narrative rehearsal. Recognizing a story and establishing a mechanism are different operations.

\section{The Ledger as a Case Study in Collection Outrunning Discrimination}

Read together, the two halves of this chapter turn an incident ledger of this kind into something more useful than a target to be individually rebutted: a large, publicly available case study in what happens when the collection of evidence proceeds without a corresponding discipline of discrimination. Entries accumulate faster than any accounting of their independence, their comparability, their underlying mechanism plurality, or their capacity to distinguish the hypothesis they are cited for from its nearest competitors.

The chain of inference under examination across the broader extinction-thesis critique developed in this Part can now be stated in a single sequence, each arrow corresponding to a distinct and separately contestable step,
\[
\text{Incidents} \xrightarrow{\text{aggregation}} \text{Evidence} \xrightarrow{\text{interpretation}} \text{Mechanism}
\]
\[
\xrightarrow{\text{extrapolation}} \text{Control} \xrightarrow{\text{prediction}} \text{Extinction},
\]
and this chapter's contribution is to show that the first arrow, from incidents to evidence, already requires an accounting of dependence, comparability, and mechanism plurality that the ledger format does not supply, while the same underlying accommodative flexibility that inflates the apparent evidence at that stage also drains the hypothesis of discriminating power at every later stage. None of this establishes that the extinction hypothesis is false. It establishes that the appropriate response to a large incident ledger is not to ask how many entries it contains, but to ask how many independent, comparable, mechanism-specific, and genuinely discriminating pieces of evidence remain once dependence, incommensurability, alarming redescription, and accommodative flexibility have all been accounted for, and what, if anything, could in principle have counted against the hypothesis the ledger is offered in support of.
